% GENERATED FILE. Edit canonical lessons, then run npm run generate:books.
% canonical-source-hash: fnv1a64:a31caae14a8f145a
% canonical-lessons: SA-C519-murtih, SA-C519-utkirati, SA-C519-silpam, SA-C519-krtih, SA-C519-nirmati

\chapter{Craft and Carving}
\label{ch:sa-craft-and-carving}
\hlchaptermodality{519}

\begin{chapteropening}
\textbf{By the end of this chapter:} \emph{I can talk about statues, carving and works of art.}

Complete the last lesson of chapter 519: I can talk about statues, carving and works of art.
\end{chapteropening}

\section[\texorpdfstring{\sk{मूर्तिः} (mūrtiḥ)}{मूर्तिः (mūrtiḥ)}]{\sk{मूर्तिः} (mūrtiḥ) — a statue}

\label{lesson:SA-C519-murtih}



\begin{quote}
Before the new one: say the Sanskrit for paints, then the Sanskrit for a portrait.
\end{quote}

\subsection*{You'll want to know: \sk{मूर्तिः}}
\textbf{\sk{मूर्तिः}} — \emph{mūrtiḥ} — \textquotedblleft{}a statue\textquotedblright{}.

\begin{quote}
\sk{मन्दिरे} \sk{एका} \sk{प्राचीना} \sk{मूर्तिः} \sk{अस्ति।}
\end{quote}

\begin{quote}
\emph{mandire ekā prācīnā mūrtiḥ asti.}
\end{quote}

\begin{quote}
There is an ancient statue in the temple.
\end{quote}

A feminine noun.

\begin{cousinweb}
\textbf{\sk{मूर्तिः}} (\emph{mūrtiḥ}) comes from the root √mūrch, whose causative MW gives as \textquotedblleft{}to make settle into a solid form\textquotedblright{}; MW sets \emph{mūrti}, any solid body or material form, an embodiment, an image, among that root's words (MW).
\end{cousinweb}

\begin{grammarlens}[title={What you've built}]
One more word for describing what you lived and felt.
\end{grammarlens}

\subsection*{Your turn}
\begin{itemize}
\raggedright
  \item \emph{Say it:} \emph{mūrtiḥ}
  \item \emph{Say it:} \emph{mūrtiḥ}, once more
  \item \emph{Say it:} \emph{mandire ekā prācīnā mūrtiḥ asti.} — There is an ancient statue in the temple.
  \item \emph{Recall:} say the Sanskrit for paints, then the Sanskrit for a portrait, then say \emph{mūrtiḥ} again
\end{itemize}

\begin{tcolorbox}[breakable,colback=teal!4,colframe=teal!35!black,title={Before you move on}]
What does \sk{मूर्तिः} mean? (\textquotedblleft{}A statue\textquotedblright{}.) Say it once more.
\end{tcolorbox}
\section[\texorpdfstring{\sk{उत्किरति} (utkirati)}{उत्किरति (utkirati)}]{\sk{उत्किरति} (utkirati) — carves}

\label{lesson:SA-C519-utkirati}



\begin{quote}
Before the new one: say the Sanskrit for a portrait, then the Sanskrit for a statue.
\end{quote}

\subsection*{You'll want to know: \sk{उत्किरति}}
\textbf{\sk{उत्किरति}} — \emph{utkirati} — \textquotedblleft{}carves\textquotedblright{}.

\begin{quote}
\sk{शिल्पी} \sk{शिलायां} \sk{पुष्पम्} \sk{उत्किरति।}
\end{quote}

\begin{quote}
\emph{śilpī śilāyāṃ puṣpam utkirati.}
\end{quote}

\begin{quote}
The sculptor carves a flower on the stone.
\end{quote}

Past (he …): \sk{सः} \sk{शिलायां} \sk{नाम} \sk{उत्कीर्णवान्}, \emph{saḥ śilāyāṃ nāma utkīrṇavān}.

\begin{cousinweb}
\textbf{\sk{उत्किरति}} (\emph{utkirati}) is built on \emph{ud}, \textquotedblleft{}up, out\textquotedblright{}, and the root √kṝ, \textquotedblleft{}to scatter\textquotedblright{}; MW gives \emph{utkṝ} as to dig out, and to engrave; its past participle is \emph{utkīrṇa} (MW).
\end{cousinweb}

\begin{grammarlens}[title={What you've built}]
One more word for describing what you lived and felt.
\end{grammarlens}

\subsection*{Your turn}
\begin{itemize}
\raggedright
  \item \emph{Say it:} \emph{utkirati}
  \item \emph{Say it:} \emph{utkirati}, once more
  \item \emph{Say it:} \emph{śilpī śilāyāṃ puṣpam utkirati.} — The sculptor carves a flower on the stone.
  \item \emph{Recall:} say the Sanskrit for a portrait, then the Sanskrit for a statue, then say \emph{utkirati} again
\end{itemize}

\begin{tcolorbox}[breakable,colback=teal!4,colframe=teal!35!black,title={Before you move on}]
What does \sk{उत्किरति} mean? (\textquotedblleft{}Carves\textquotedblright{}.) Say it once more.
\end{tcolorbox}
\section[\texorpdfstring{\sk{शिल्पम्} (śilpam)}{शिल्पम् (śilpam)}]{\sk{शिल्पम्} (śilpam) — a craft}

\label{lesson:SA-C519-silpam}



\begin{quote}
Before the new one: say the Sanskrit for a statue, then the Sanskrit for carves.
\end{quote}

\subsection*{You'll want to know: \sk{शिल्पम्}}
\textbf{\sk{शिल्पम्}} — \emph{śilpam} — \textquotedblleft{}a craft\textquotedblright{}.

\begin{quote}
\sk{शिल्पं} \sk{कठिनं} \sk{परन्तु} \sk{सुन्दरम्} \sk{अस्ति।}
\end{quote}

\begin{quote}
\emph{śilpaṃ kaṭhinaṃ parantu sundaram asti.}
\end{quote}

\begin{quote}
A craft is hard but beautiful.
\end{quote}

A neuter noun.

\begin{cousinweb}
\textbf{\sk{शिल्पम्}} (\emph{śilpam}) is the word MW gives for decoration, ornament, artistic work, and MW itself calls it of doubtful derivation, so its origin is uncertain (MW).
\end{cousinweb}

\begin{grammarlens}[title={What you've built}]
One more word for describing what you lived and felt.
\end{grammarlens}

\subsection*{Your turn}
\begin{itemize}
\raggedright
  \item \emph{Say it:} \emph{śilpam}
  \item \emph{Say it:} \emph{śilpam}, once more
  \item \emph{Say it:} \emph{śilpaṃ kaṭhinaṃ parantu sundaram asti.} — A craft is hard but beautiful.
  \item \emph{Recall:} say the Sanskrit for a statue, then the Sanskrit for carves, then say \emph{śilpam} again
\end{itemize}

\begin{tcolorbox}[breakable,colback=teal!4,colframe=teal!35!black,title={Before you move on}]
What does \sk{शिल्पम्} mean? (\textquotedblleft{}A craft\textquotedblright{}.) Say it once more.
\end{tcolorbox}
\section[\texorpdfstring{\sk{कृतिः} (kṛtiḥ)}{कृतिः (kṛtiḥ)}]{\sk{कृतिः} (kṛtiḥ) — a creation; a work}

\label{lesson:SA-C519-krtih}



\begin{quote}
Before the new one: say the Sanskrit for carves, then the Sanskrit for a craft.
\end{quote}

\subsection*{You'll want to know: \sk{कृतिः}}
\textbf{\sk{कृतिः}} — \emph{kṛtiḥ} — \textquotedblleft{}a creation; a work\textquotedblright{}.

\begin{quote}
\sk{एषा} \sk{मूर्तिः} \sk{मम} \sk{पितुः} \sk{कृतिः} \sk{अस्ति।}
\end{quote}

\begin{quote}
\emph{eṣā mūrtiḥ mama pituḥ kṛtiḥ asti.}
\end{quote}

\begin{quote}
This statue is my father's work.
\end{quote}

A feminine noun.

\begin{cousinweb}
\textbf{\sk{कृतिः}} (\emph{kṛtiḥ}) comes from the root √kṛ, \textquotedblleft{}to make, to do\textquotedblright{}, with the ending \emph{-ti}; MW gives \emph{kṛti} as the act of doing, a creation, a work (MW).
\end{cousinweb}

\begin{grammarlens}[title={What you've built}]
One more word for describing what you lived and felt.
\end{grammarlens}

\subsection*{Your turn}
\begin{itemize}
\raggedright
  \item \emph{Say it:} \emph{kṛtiḥ}
  \item \emph{Say it:} \emph{kṛtiḥ}, once more
  \item \emph{Say it:} \emph{eṣā mūrtiḥ mama pituḥ kṛtiḥ asti.} — This statue is my father's work.
  \item \emph{Recall:} say the Sanskrit for carves, then the Sanskrit for a craft, then say \emph{kṛtiḥ} again
\end{itemize}

\begin{tcolorbox}[breakable,colback=teal!4,colframe=teal!35!black,title={Before you move on}]
What does \sk{कृतिः} mean? (\textquotedblleft{}A creation; a work\textquotedblright{}.) Say it once more.
\end{tcolorbox}
\section[\texorpdfstring{\sk{निर्माति} (nirmāti)}{निर्माति (nirmāti)}]{\sk{निर्माति} (nirmāti) — fashions; builds}

\label{lesson:SA-C519-nirmati}



\begin{quote}
Before the new one: say the Sanskrit for a craft, then the Sanskrit for a creation.
\end{quote}

\subsection*{You'll want to know: \sk{निर्माति}}
\textbf{\sk{निर्माति}} — \emph{nirmāti} — \textquotedblleft{}fashions; builds\textquotedblright{}.

\begin{quote}
\sk{शिल्पी} \sk{शिलायाः} \sk{गजं} \sk{निर्माति।}
\end{quote}

\begin{quote}
\emph{śilpī śilāyāḥ gajaṃ nirmāti.}
\end{quote}

\begin{quote}
The craftsman fashions an elephant out of stone.
\end{quote}

Past (he …): \sk{सः} \sk{शिलायाः} \sk{गजं} \sk{निर्मितवान्}, \emph{saḥ śilāyāḥ gajaṃ nirmitavān}.

\begin{cousinweb}
\textbf{\sk{निर्माति}} (\emph{nirmāti}) is built on \emph{nis}, \textquotedblleft{}out\textquotedblright{}, and the root √mā, \textquotedblleft{}to measure, to form\textquotedblright{}; MW gives \emph{nirmā} as to build, to make out of a material, to create; its past participle is \emph{nirmita} (MW).
\end{cousinweb}

\begin{grammarlens}[title={What you've built}]
One more word for describing what you lived and felt.
\end{grammarlens}

\subsection*{Your turn}
\begin{itemize}
\raggedright
  \item \emph{Say it:} \emph{nirmāti}
  \item \emph{Say it:} \emph{nirmāti}, once more
  \item \emph{Say it:} \emph{śilpī śilāyāḥ gajaṃ nirmāti.} — The craftsman fashions an elephant out of stone.
  \item \emph{Recall:} say the Sanskrit for a craft, then the Sanskrit for a creation, then say \emph{nirmāti} again
\end{itemize}

\begin{tcolorbox}[breakable,colback=teal!4,colframe=teal!35!black,title={Before you move on}]
What does \sk{निर्माति} mean? (\textquotedblleft{}Fashions; builds\textquotedblright{}.) Say it once more.
\end{tcolorbox}
